\section{Existing ML benchmarks for distribution shifts}\label{sec:related}
Distribution shifts have been a longstanding problem in the ML research community \citep{hand2006classifier,quinonero2009dataset}. Earlier work studied shifts in datasets for tasks including
part-of-speech tagging \citep{marcus93treebank},
sentiment analysis \citep{blitzer2007biographies},
land cover classification \citep{bruzzone2009domain},
object recognition \citep{saenko2010adapting},
and flow cytometry \citep{blanchard2011generalizing}.
However, these datasets are not as widely used today, in part because they tend to be much smaller than modern datasets.

Instead, many recent papers have focused on object recognition datasets with shifts induced by synthetic transformations, such as
ImageNet-C \citep{hendrycks2019benchmarking}, which corrupts images with noise;
the Backgrounds Challenge \citep{xiao2020noise} and Waterbirds \citep{sagawa2020group}, which alter image backgrounds;
or Colored MNIST \citep{arjovsky2019invariant}, which changes the colors of MNIST digits.
It is also common to use data splits or combinations of disparate datasets to induce shifts, such as generalizing to photos solely from cartoons and other stylized images in PACS \citep{li2017deeper}; generalizing to objects at different scales solely from a single scale in DeepFashion Remixed \citep{hendrycks2020many}; or using training and test sets with disjoint subclasses in BREEDS \citep{santurkar2020breeds} and similar datasets \citep{hendrycks2019benchmarking}. While our treatment here is necessarily brief, we discuss other similar datasets in \refapp{app_related}.

These existing benchmarks are useful and important testbeds for method development.
As they typically target well-defined and isolated shifts,
they facilitate clean analysis and controlled experimentation,
e.g., studying the effect of backgrounds on image classification \citep{xiao2020noise}, or showing that training with added Gaussian blur improves performance on real-world blurry images \citep{hendrycks2020many}.
Moreover, by studying how off-the-shelf models trained on standard datasets like ImageNet perform on different test datasets, we can better understand the robustness of these widely-used models
\citep{geirhos2018generalisation, recht2019doimagenet, hendrycks2019benchmarking, taori2020measuring, djolonga2020robustness, hendrycks2020many}.

However, as we discussed in the introduction, robustness to these synthetic shifts need not transfer to the kinds of shifts that arise in real-world deployments \citep{taori2020measuring,djolonga2020robustness,xie2020innout},
and it is thus challenging to develop and evaluate methods for training models that are robust to real-world shifts on these datasets alone.
With WILDS, we seek to complement existing benchmarks by curating datasets that reflect natural distribution shifts across a diverse set of data modalities and application.
